\chapter{Integrated Review and Cache Simulation Concepts}\label{ch:auto-006}


\section{Chapter Purpose}\label{sec:auto-046}
This chapter integrates the concepts from Chapters 1--4 into a single combined overview.


\section{Integrated Themes from Chapters 1 and 2}\label{sec:auto-047}
Core concept clusters include classifying optimizations as parallelism, speculation, or both; distinguishing throughput from runtime and selecting the correct metric; interpreting $\text{CPU time}=\text{IC}\times\text{CPI}\times\text{CT}$ as a tradeoff framework; applying Amdahl's law for fixed-size speedup limits; interpreting Gustafson's law for scaled-size behavior; and using arithmetic means for times versus harmonic means for rates.

\paragraph{Worked check}\label{par:auto-026}
If an optimization accelerates 30\% of runtime by $2\times$, total speedup is
\begin{equation}\label{eq:auto-045}
    S=\frac{1}{(1-0.3)+\frac{0.3}{2}}=\frac{1}{0.85}\approx 1.176.
\end{equation}
The unaccelerated fraction remains in the denominator and bounds overall speedup.

\paragraph{Worked check: maximum possible speedup}\label{par:auto-026a}
If a mechanism accelerates 25\% of the original benchmark runtime and that portion is made arbitrarily fast, the maximum total speedup is
\begin{equation}\label{eq:auto-045a}
    S_{\max}=\frac{1}{1-0.25}=\frac{4}{3}\approx 1.33.
\end{equation}
This is the cleanest way to read the \emph{maximum} speedup of a partial optimization.


\section{Integrated Themes from Chapter 3}\label{sec:auto-048}
Core concept clusters include computing sets, ways, and tag/index/offset widths from $(C,B,S)$; interpreting compulsory, conflict, and capacity misses; analyzing replacement and write-policy behavior; applying $AAT=HT+MR\cdot MP$; and relating cache size, block size, and associativity to hit time, miss ratio, and miss penalty.

\paragraph{Worked check}\label{par:auto-027}
For a cache with $HT=1.8$ ns, $MR=0.06$, and $MP=28$ ns (Eq.~\eqref{eq:auto-046}):
\begin{equation}\label{eq:auto-046}
    AAT = 1.8 + 0.06\cdot 28 = 3.48\ \text{ns}.
\end{equation}

\paragraph{Worked check: repeated trace reasoning}\label{par:auto-027a}
Consider a direct-mapped L1 with $(C,B,S)=(8,4,0)$, write-back/write-allocate, and the repeated access sequence
\begin{equation}\label{eq:auto-046a}
    R\ 0x41F1,\ R\ 0xB54A,\ W\ 0x4204,\ W\ 0xCA3B,\ R\ 0x420E,\ W\ 0xB537.
\end{equation}
With $B=4$, the six accesses touch the five blocks $\{0x41F,0xB54,0x420,0xCA3,0xB53\}$ because $0x4204$ and $0x420E$ lie in the same block.
The first iteration therefore has 5 compulsory misses.
The key structural fact is that blocks $0xCA3$ and $0xB53$ map to the same direct-mapped set, so after warm-up the sequence continues to suffer two conflict misses per iteration.
Over iterations 2--10, that produces 18 additional misses.
Those 18 misses are the non-compulsory misses for the full run, and because the conflicting writes repeatedly evict dirty lines, the full sequence produces 19 writebacks if dirty lines are \emph{not} flushed after the final access.

\paragraph{Victim cache and +1 prefetch contrast}\label{par:auto-027b}
That same trace separates two common remedies cleanly.
A small victim cache can eliminate the recurring two-block ping-pong at one set; after warm-up, the combined L1+victim-cache miss count can drop to zero because each conflicting block is recovered from the victim structure instead of DRAM\@.
By contrast, a $+1$ prefetcher is not automatically a cure.
In this trace it can prefetch block $0x420$ early enough to remove one first-iteration demand miss, but it does not remove the recurring conflict between $0xCA3$ and $0xB53$.
This is the key design lesson of the trace: next-block prefetching and victim caching address different miss mechanisms.


\section{Worked Checks for Chapter 4}\label{sec:auto-049}
Core concept clusters from Chapter 4 include reading timing diagrams to infer hazards and hardware support, distinguishing separate instruction and data-memory paths from shared-memory structural hazards, identifying which bypass paths are present, checking whether the register file supports same-cycle write-then-read behavior, and simulating branch predictors such as always taken, one-bit and Smith-counter schemes, GShare, GSelect, Yeh/Patt, perceptrons, and hybrid tournament predictors.

\subsection*{Pipeline Speedup}\label{subsec:auto-012}
Assume a 5-stage design where each stage has 250 ps logic delay and pipeline-register overhead is 30 ps.
Unpipelined instruction time is $5\times250=1250$ ps.
Pipelined clock is $250+30=280$ ps.
Ideal long-run speedup is (Eq.~\eqref{eq:auto-047}):
\begin{equation}\label{eq:auto-047}
    S\approx\frac{1250}{280}\approx 4.46.
\end{equation}
This is below 5 because register overhead is non-zero.

\subsection*{CPI with Branch Stalls}\label{subsec:auto-013}
If branch frequency is 18\% and each branch contributes an average of 1.6 stall cycles (Eq.~\eqref{eq:auto-048}):
\begin{equation}\label{eq:auto-048}
    CPI = 1 + 0.18\times1.6 = 1.288.
\end{equation}
This illustrates how even moderate branch penalties raise CPI above the ideal value of 1.

\subsection*{Software Branch Handling}\label{subsec:auto-013a}
When a branch has delay slots, every instruction moved into those slots must be safe regardless of whether the branch is taken.
When a branch instead has squash slots under an ``assume taken'' policy, the compiler may place instructions from the predicted taken path there, because wrong-path entries will be canceled.
So the safe rule is simple:
\begin{enumerate}
    \item fill delay slots only with path-independent instructions,
    \item fill squash slots from the predicted path when profitable, and
    \item leave \texttt{NOP}s when correctness would otherwise be uncertain.
\end{enumerate}


\section{Cache Simulation Conceptual Model}\label{sec:auto-050}
The cache simulation framework models hierarchy behavior at the level of addresses, tags, and policy events rather than payload data values.
A common parameterization is (Eq.~\eqref{eq:auto-049}):
\begin{equation}\label{eq:auto-049}
    (C_1,B,S_1,C_2,S_2,\{\text{MIP,LIP}\},\{\text{None,+1,Markov,Hybrid}\},r).
\end{equation}
In this notation, $C$ controls capacity ($2^C$ bytes), $B$ controls block size ($2^B$ bytes), and $S$ controls associativity ($2^S$ ways).
Here \textit{most recently used insertion policy (MIP)} and \textit{least recently used insertion policy (LIP)} are common insertion-policy options.


\section{Hierarchy Semantics in the Model}\label{sec:auto-051}
\textbf{L1/L2 abstraction:} L1 accesses are modeled independently from L2 internals, and L2 requests are modeled as arriving from an upper-level requester.
\textbf{Write-strategy contrast:} \textit{write-back, write-allocate (WBWA)} behavior in L1 emphasizes dirty-state tracking and deferred propagation, while \textit{write-through, write-no-allocate (WTWNA)} behavior in L2 emphasizes immediate lower-level visibility and no allocation on write misses.

\textbf{Miss-servicing order effects:} ordering choices in read-miss repair and writeback events influence replacement-state evolution and therefore downstream statistics.
\textbf{Disabled-level interpretation:} a disabled L2 can be represented as an always-miss transit layer with zero L2 hit time in AAT terms.


\section{Prefetching Concepts}\label{sec:auto-052}
\textbf{$+1$ prefetching:} predicts the next sequential block ($X+1$) after a demand miss on block $X$.

\textbf{Markov prefetching:}
represents miss-stream transitions as a finite-state table, associates each source miss block with a bounded set of likely successor blocks and frequencies, and selects a prediction from the highest-frequency successor under a deterministic tie-break rule.

\textbf{Hybrid prefetching:} combines correlation-based prediction with sequential fallback when correlation information is absent.

\textbf{Prefetch usefulness tracking:}
a prefetch metadata bit distinguishes demand-filled lines from prefetched lines.
First demand use of a prefetched line contributes to prefetch-hit interpretation, while eviction before demand use contributes to prefetch-miss interpretation.


\section{Statistics and AAT Interpretation}\label{sec:auto-053}
Typical hierarchy statistics include access counts, hit/miss counts, hit/miss ratios, writeback counts, and prefetch effectiveness counters.
In two-level analysis (Eq.~\eqref{eq:auto-050}):
\begin{equation}\label{eq:auto-050}
    AAT_{L1}=HT_{L1}+MR_{L1}\cdot MP_{L1},\qquad MP_{L1}=AAT_{L2}.
\end{equation}
\begin{equation}\label{eq:auto-051}
    AAT_{L2}=HT_{L2}+MR_{L2,read}\cdot DRAM\_TIME.
\end{equation}
\begin{equation}\label{eq:auto-052}
    DRAM\_TIME = DRAM\_AT + DRAM\_AT\_PER\_WORD\cdot WORDS\_PER\_BLOCK.
\end{equation}
This form emphasizes that higher-level penalty is the lower-level average access behavior, not a fixed constant.


\section{Conceptual Tradeoffs Across Parameters}\label{sec:auto-054}
Increasing $C$ generally reduces capacity misses but tends to raise hit time.
Increasing $S$ reduces conflict misses but increases lookup complexity.
Increasing $B$ can exploit spatial locality but may increase pollution and miss penalty.
More aggressive prefetching can reduce compulsory misses yet increase bandwidth pressure and useless fills.


\section{Chapter 5 Summary}\label{sec:auto-055}
Integrated themes, pipeline analysis, and cache simulation models reinforce the same central idea: architecture performance is governed by measurable tradeoffs among latency, throughput, locality, and prediction quality.
The simulation concepts provide a concrete way to observe those tradeoffs in controlled workloads.
That same simulator mindset extends naturally to branch-prediction accuracy, fetched-versus-retired instruction counts, queue pressure, and reorder-buffer pressure, which are interpreted more fully in Chapter 6.
A grouped formula index for the full text is provided in Appendix~\ref{ch:formula-compendium}.
